\documentclass[12pt,a4paper]{article}

% ── Packages ──────────────────────────────────────────────────────────────────
\usepackage[utf8]{inputenc}
\usepackage[T1]{fontenc}
\usepackage{lmodern}
\usepackage[margin=1in]{geometry}
\usepackage{amsmath, amssymb, amsthm}
\usepackage{graphicx}
\usepackage{hyperref}
\usepackage{booktabs}
\usepackage{enumitem}
\usepackage{xcolor}
\usepackage{algorithm}
\usepackage{algpseudocode}
\usepackage{natbib}
\usepackage{tikz}
\usetikzlibrary{arrows.meta, positioning, shapes.geometric}

% ── Metadata ──────────────────────────────────────────────────────────────────
\title{%
  \textbf{Applied Cognitive Science $\times$ Reinforcement Learning $\times$ LLMs}\\[0.4em]
  \large Towards Principled Enterprise-Scale AI Agents\\[0.2em]
  \normalsize Thursday Learning Hours — Session 1
}
\author{Prabakaran C}
\date{\today}

% ── Theorem environments ───────────────────────────────────────────────────────
\newtheorem{definition}{Definition}
\newtheorem{theorem}{Theorem}
\newtheorem{proposition}{Proposition}
\newtheorem{remark}{Remark}

% ──────────────────────────────────────────────────────────────────────────────
\begin{document}

\maketitle
\tableofcontents
\newpage

% ══════════════════════════════════════════════════════════════════════════════
\section{Introduction}
% ══════════════════════════════════════════════════════════════════════════════

Large Language Models (LLMs) have reshaped the landscape of enterprise AI. Yet most deployed
``agentic'' systems remain shallow: a model wrapped in a prompt template, a retrieval pipeline,
and an orchestration harness. While effective for narrow tasks, such scaffolding fails in
settings that require long-horizon planning, adaptive decision-making under uncertainty, and
genuine tool-use strategy.

This work argues that two mature fields — \textbf{Cognitive Science} and
\textbf{Reinforcement Learning (RL)} — provide the missing theoretical and practical
foundations for building robust, enterprise-grade AI agents. We survey the relevant literature,
identify open problems, and present a series of experiments that demonstrate concrete gains
over scaffold-only baselines.

% ══════════════════════════════════════════════════════════════════════════════
\section{Background}
% ══════════════════════════════════════════════════════════════════════════════

\subsection{Cognitive Architectures}

Cognitive architectures model the functional organisation of an intelligent mind. Influential
examples include:

\begin{itemize}[noitemsep]
  \item \textbf{ACT-R} \citep{anderson2004integrated} — procedural and declarative memory modules
        with activation-based retrieval.
  \item \textbf{SOAR} \citep{laird2012soar} — production-rule system with an explicit problem space.
  \item \textbf{Global Workspace Theory (GWT)} \citep{baars1988cognitive} — a ``broadcast'' model of
        conscious attention.
  \item \textbf{Predictive Processing} \citep{clark2016surfing} — the brain as a hierarchical
        prediction machine minimising free energy.
\end{itemize}

\subsection{Reinforcement Learning Fundamentals}

We model an agent's interaction with the world as a Markov Decision Process
$\mathcal{M} = \langle \mathcal{S}, \mathcal{A}, P, R, \gamma \rangle$ where
$\mathcal{S}$ is the state space, $\mathcal{A}$ the action space,
$P: \mathcal{S} \times \mathcal{A} \to \Delta(\mathcal{S})$ the transition kernel,
$R: \mathcal{S} \times \mathcal{A} \to \mathbb{R}$ the reward function, and
$\gamma \in [0,1)$ a discount factor.

The agent seeks a policy $\pi^*$ that maximises the expected discounted return:
\begin{equation}
  \pi^* = \arg\max_\pi \; \mathbb{E}_\pi \left[ \sum_{t=0}^{\infty} \gamma^t R(s_t, a_t) \right].
\end{equation}

\subsection{RLHF and Post-Training Alignment}

Reinforcement Learning from Human Feedback (RLHF) \citep{christiano2017deep} finetunes LLMs
using a learned reward model trained on human preference pairs. Extensions such as
RLAIF \citep{bai2022constitutional} and GRPO \citep{shao2024deepseekmath} reduce reliance on
human annotators and improve sample efficiency.

% ══════════════════════════════════════════════════════════════════════════════
\section{Cognitive Science $\rightarrow$ Agent Design Principles}
% ══════════════════════════════════════════════════════════════════════════════

\subsection{Working Memory and Context Management}

GWT's global workspace maps naturally onto the LLM context window. We propose a
\emph{selective broadcast} mechanism: only high-salience state representations are
admitted into context, reducing noise and improving long-horizon coherence.

\subsection{Episodic and Semantic Memory}

Drawing on the declarative memory distinction in ACT-R:

\begin{itemize}[noitemsep]
  \item \textbf{Episodic store}: vector database of past interaction traces, retrieved by
        similarity to the current query.
  \item \textbf{Semantic store}: structured knowledge graph of domain facts, queried
        symbolically.
\end{itemize}

\subsection{Metacognition and Uncertainty Estimation}

Predictive processing suggests agents should maintain calibrated beliefs about their own
competence. We implement this as a lightweight \emph{confidence head} attached to the policy
network that gates whether the LLM should act, query memory, or ask for human input.

% ══════════════════════════════════════════════════════════════════════════════
\section{RL for Enterprise Agentic Systems}
% ══════════════════════════════════════════════════════════════════════════════

\subsection{Reward Shaping for Tool Use}

Let $\mathcal{T} = \{t_1, \ldots, t_K\}$ be the tool set. We define a shaped reward:
\begin{equation}
  \tilde{R}(s, a) = R(s, a) + \phi(s') - \phi(s),
\end{equation}
where $\phi: \mathcal{S} \to \mathbb{R}$ is a potential function that assigns credit for
intermediate task milestones (e.g., successfully fetching a required data source).

\subsection{Hierarchical RL for Long-Horizon Tasks}

We decompose tasks into a two-level hierarchy:
\begin{itemize}[noitemsep]
  \item \textbf{High-level policy} $\pi^H$: selects sub-goals $g \in \mathcal{G}$.
  \item \textbf{Low-level policy} $\pi^L(a \mid s, g)$: executes primitive actions to reach $g$.
\end{itemize}
This mirrors the goal/sub-goal decomposition observed in human problem solving \citep{miller1960plans}.

\subsection{Model-Based RL and World Models}

A world model $\hat{P}(s' \mid s, a)$ enables the agent to plan in imagination before
committing to real actions — directly analogous to mental simulation in cognitive science.
We evaluate Dreamer-V3 \citep{hafner2023mastering} as a backbone for enterprise workflow agents.

% ══════════════════════════════════════════════════════════════════════════════
\section{Experiments}
% ══════════════════════════════════════════════════════════════════════════════

\subsection{Setup}

All experiments use a GPT-4o or open-weight Llama-3 backbone.
Baselines: (i) vanilla ReAct \citep{yao2023react}, (ii) LangChain agent harness.
Our systems additionally incorporate the cognitive and RL modules described above.

\subsection{Benchmarks}

\begin{table}[h!]
\centering
\caption{Benchmark task categories}
\begin{tabular}{@{}lll@{}}
\toprule
\textbf{Task Type} & \textbf{Dataset} & \textbf{Metric} \\
\midrule
Multi-step QA        & HotpotQA                   & F1 / EM \\
Tool-use planning    & ToolBench                  & Success Rate \\
Code generation      & SWE-bench                  & \% resolved \\
Enterprise workflow  & Custom (internal)          & Task Completion \\
\bottomrule
\end{tabular}
\end{table}

\subsection{Results}

\textit{[Results to be populated after experiments — see \texttt{experiments/} folder.]}

% ══════════════════════════════════════════════════════════════════════════════
\section{Discussion}
% ══════════════════════════════════════════════════════════════════════════════

\subsection{Limitations}

\begin{itemize}[noitemsep]
  \item Reward design remains the key bottleneck; bad reward functions induce reward hacking.
  \item World models are expensive to train and may hallucinate dynamics at distribution shift.
  \item Cognitive architecture mappings are loose analogies; neuroscientific grounding is partial.
\end{itemize}

\subsection{Future Directions}

\begin{enumerate}[noitemsep]
  \item Tighter integration of spiking neural network models of timing and attention.
  \item Multi-agent RL with cognitive role specialisation (planner / executor / critic).
  \item Continual RL for agents deployed over months in live enterprise settings.
\end{enumerate}

% ══════════════════════════════════════════════════════════════════════════════
\section{Conclusion}
% ══════════════════════════════════════════════════════════════════════════════

Scaffolding is a starting point, not a destination. By grounding LLM agent design in
cognitive science and reinforcement learning, we can build systems that learn, plan, and
adapt — the hallmarks of genuine intelligence. This sprint lays the theoretical and
experimental foundation for that programme.

% ── Bibliography ──────────────────────────────────────────────────────────────
\newpage
\bibliographystyle{plainnat}
\bibliography{references}

% Inline bibliography placeholder (replace with .bib file)
\begin{thebibliography}{99}

\bibitem{anderson2004integrated}
Anderson, J.~R. et al. (2004).
\newblock An integrated theory of the mind.
\newblock \textit{Psychological Review}, 111(4), 1036.

\bibitem{laird2012soar}
Laird, J.~E. (2012).
\newblock \textit{The Soar Cognitive Architecture}.
\newblock MIT Press.

\bibitem{baars1988cognitive}
Baars, B.~J. (1988).
\newblock \textit{A Cognitive Theory of Consciousness}.
\newblock Cambridge University Press.

\bibitem{clark2016surfing}
Clark, A. (2016).
\newblock \textit{Surfing Uncertainty: Prediction, Action, and the Embodied Mind}.
\newblock Oxford University Press.

\bibitem{christiano2017deep}
Christiano, P. et al. (2017).
\newblock Deep reinforcement learning from human preferences.
\newblock \textit{NeurIPS}.

\bibitem{bai2022constitutional}
Bai, Y. et al. (2022).
\newblock Constitutional AI: Harmlessness from AI Feedback.
\newblock \textit{arXiv:2212.08073}.

\bibitem{shao2024deepseekmath}
Shao, Z. et al. (2024).
\newblock DeepSeekMath: Pushing the Limits of Mathematical Reasoning in Open Language Models.
\newblock \textit{arXiv:2402.03300}.

\bibitem{miller1960plans}
Miller, G.~A., Galanter, E., \& Pribram, K.~H. (1960).
\newblock \textit{Plans and the Structure of Behavior}.
\newblock Henry Holt.

\bibitem{hafner2023mastering}
Hafner, D. et al. (2023).
\newblock Mastering Diverse Domains through World Models.
\newblock \textit{arXiv:2301.04104}.

\bibitem{yao2023react}
Yao, S. et al. (2023).
\newblock ReAct: Synergizing Reasoning and Acting in Language Models.
\newblock \textit{ICLR 2023}.

\end{thebibliography}

\end{document}
